% Answers to Chapter 6, translated from sv/back/facit/06.tex. Then the answers to practice test 1.

\facitavsnitt{c6}{Chapter 6}{Quadratic equations}
\begin{facit}

\svar{1.1} \sv{a} $10$ \sv{b} $4 \cdot 10^6$ \sv{c} $2 \times 10^3 = 2\,000$

\svar{1.2} \sv{a} $I \approx 7.02 \times 10^{-1}\,\text{mA} = 0.702\,\text{mA}$
\sv{b} $P \approx 2.32\,\text{mW}$

\svar{1.3} \sv{a} $12$ \sv{b} $0.07$ \sv{c} $20\,\text{V}$

\svar{1.4} \sv{a} $0100\,0100_2 = \mathrm{44}_{16}$ \sv{b} $0000\,0100_2 = \mathrm{04}_{16}$
\sv{c} $x = 0100\,0000_2 = \mathrm{40}_{16}$, and \code{PORTD} becomes
$0100\,0100_2 = \mathrm{44}_{16}$.
\sv{d} $R = 300\,\Omega$ gives exactly $10\,\text{mA}$; the nearest E12 values are $330\,\Omega$
($I \approx 9.1\,\text{mA}$) and $270\,\Omega$ ($I \approx 11.1\,\text{mA}$).
\sv{e} About $45\,\text{mW}$ per circuit with $330\,\Omega$ (about $56\,\text{mW}$ with
$270\,\Omega$).

\svar{2.1} \sv{a} $I = \pm 0.4\,\text{A}$
\sv{b} The negative root only says that the current flows the other way; the power $RI^2$ is the
same, so the magnitude of the current is $I = 0.4\,\text{A}$.

\svar{2.2} \sv{a} $x_1 = 3$, $x_2 = 2$ \sv{b} $x_1 = 2$, $x_2 = -5$
\sv{c} $x_1 = x_2 = 3$ (a double root)

\svar{2.3} \sv{a} $x_1 = 3$, $x_2 = 0.5$ \sv{b} $x_1 = \frac{1}{3}$, $x_2 = -2$

\svar{2.4} \sv{a} $R_1^2 - 13R_1 + 36 = 0$ \sv{b} $9\,\Omega$ and $4\,\Omega$

\svar{3.1} \sv{a} $x = \pm 7$ \sv{b} $x = \pm 3$ \sv{c} $x = \pm 5$
\sv{d} No real solutions, since $x^2 = -4$. \sv{e} $x = \pm 6$

\svar{3.2} \sv{a} $x_1 = 0$, $x_2 = 7$ \sv{b} $x = \pm 4$ \sv{c} $x_1 = -2$, $x_2 = -3$
\sv{d} $x_1 = 0$, $x_2 = 4$ \sv{e} $x_1 = x_2 = 2$ (a double root)

\svar{3.3} \sv{a} $x_1 = 2$, $x_2 = -4$ \sv{b} $x_1 = 6$, $x_2 = 2$ \sv{c} $x_1 = 3$, $x_2 = -4$
\sv{d} $x_1 = 5$, $x_2 = -3$

\svar{3.4} \sv{a} $x_1 = 3$, $x_2 = 2$ \sv{b} $x_1 = 1$, $x_2 = -5$ \sv{c} $x_1 = 3$, $x_2 = 1$

\svar{3.5} The expression under the root sign, $\left(\frac{p}{2}\right)^2 - q$, is:
\sv{a} $4 > 0$, two real roots \sv{b} $0$, a double root
\sv{c} $-2.75 < 0$, no real roots \sv{d} $0$, a double root
\sv{e} $\frac{4}{9} > 0$, two real roots

\svar{3.6} \sv{a} $x_1 = 1$, $x_2 = -5$ \sv{b} $x_1 = 7$, $x_2 = 3$
\sv{c} $x = -3 \pm \sqrt{5}$, so $x_1 \approx -0.76$ and $x_2 \approx -5.24$

\svar{3.7} \sv{a} $0.5\,\text{A}$ \sv{b} $30\,\text{mA}$ \sv{c} About $33\,\text{mA}$
\sv{d} The negative root only means the opposite direction of current; the power is the same
because $I$ is squared, and it is the magnitude of the current that is asked for.

\svar{3.8} \sv{a} $10\,\text{V}$ \sv{b} About $15.8\,\text{V}$ \sv{c} $52.9\,\Omega$

\svar{3.9} \sv{a} $12\,\Omega$ and $8\,\Omega$ \sv{b} $8\,\Omega$ and $7\,\Omega$
\sv{c} No. For $R_1^2 - 10R_1 + 30 = 0$ the expression under the root sign is
$25 - 30 = -5 < 0$, so there are no real roots.

\svar{3.10} \sv{a} Product over sum gives $R_1 R_2 = (R_1 \parallell R_2)(R_1 + R_2)
= 2.4 \times 10 = 24\,\Omega^2$.
\sv{b} $R_1^2 - 10R_1 + 24 = 0$ gives $6\,\Omega$ and $4\,\Omega$.
\sv{c} $6 + 4 = 10\,\Omega$ and $\frac{6 \times 4}{6 + 4} = 2.4\,\Omega$.

\svar{3.11} \sv{a} $P = R\left(\frac{12}{4 + R}\right)^2 = 8$ gives $18R = (4 + R)^2$, so
$R^2 - 10R + 16 = 0$.
\sv{b} $R = 2\,\Omega$ or $R = 8\,\Omega$
\sv{c} $R = 2\,\Omega$: $I = 2\,\text{A}$ and $P = 8\,\text{W}$. $R = 8\,\Omega$: $I = 1\,\text{A}$
and $P = 8\,\text{W}$.
\sv{d} $I = 1.5\,\text{A}$ and $P = 9\,\text{W}$, more than the two solutions' $8\,\text{W}$: the
power in $R$ is greatest when $R = R_1$ (maximum power transfer).

\svar{3.12} \sv{a} False: the expression under the root sign in the PQ-formula can give two, one
or no real roots.
\sv{b} False: $x = -4$ is also a solution.
\sv{c} False: the equation must first be normalised, $x^2 + 2x - 3 = 0$.
\sv{d} True: the square-root term becomes zero, and both roots become $-\frac{p}{2}$.
\sv{e} True: that is the zero-product rule.
\sv{f} True: $x^2 = -4$, and no real square is negative.

\svar[Test yourself]{T} \sv{1} $2x(x - 4) = 0$, so $x_1 = 0$ and $x_2 = 4$.
\sv{2} $I = 0.5$ (amperes). The negative root is rejected: it only indicates the opposite
direction of current, and the power $RI^2$ is the same.

\end{facit}

\facitavsnitt{od1}{Practice test 1}{Arithmetic, algebra, powers, number systems and equations}
\begin{facit}

\svar{1} $R_2 \approx 8\,\text{k}\Omega$ (unrounded $8.18\,\text{k}\Omega$)

\svar{2} $c(c^2 + 1) = c^3 + c$

\svar{3} $x = -\frac{17}{7}$, $y = -\frac{19}{7}$

\svar{4} $1001\,1100_2 = 156$

\svar{5} $P \approx 2.9 \cdot 10^{-2}\,\text{W}$ (about $29\,\text{mW}$)

\svar{6} $x_1 \approx 0.21$, $x_2 \approx -27.54$

\end{facit}
